All constants below depend only on \(q,L,c_-,r_0\).  We first construct one estimator that works for every law in the class.  For an arm \(t\) and boundary point \(x\), write
\[
W_{i,t,x}(h)=\mathbf1\{T_i=t\}K_\triangle(D_{P,x,i}/h),
\qquad
G_{t,x}(h)=\frac1{nh^2}\sum_{i=1}^nW_{i,t,x}(h).
\]
The local-constant estimate is the weighted outcome average when \(G_{t,x}(h)>0\), and is zero otherwise.  It is a measurable function of \(U_{n,P}(x)\), so its arm difference belongs to \(\mathcal T_n^{\mathrm{dist}}\).  By the stability lemma,
\[
\inf_{P,t,x}\mathbb E_PG_{t,x}(h)\geq\kappa_0
\]
for \(h\leq2r_0\).  By the upper-mass lemma, a kernel summand is nonzero with probability at most \(\pi Lh^2\).

We record the uniform stochastic calculation.  Conditional sub-Gaussianity, boundedness of the kernel, and the preceding support probability imply the Bernstein moment bounds
\[
\mathbb E_P|W_{i,t,x}(h)\{Y_i-\mu_{t,P}(X_i)\}|^k
\leq k!C^k h^2,
\qquad k\geq2,
\]
uniformly in \(P,t,x\); the same bounds without the outcome factor hold for centered Gram summands.  Exponential Markov optimization therefore gives, for \(0<u\leq1\),
\[
\Pr_P\!\left(\left|G_{t,x}(h)-\mathbb E_PG_{t,x}(h)\right|>u\right)
+\Pr_P\!\left(\left|\frac1{nh^2}\sum_{i=1}^nW_{i,t,x}(h)
\{Y_i-\mu_{t,P}(X_i)\}\right|>u\right)
\leq C e^{-cnh^2u^2}.
\]

Choose a maximal \(n^{-3}\)-separated subset of \(\mathcal B_P\).  Since \(\mathcal B_P\subseteq[-L,L]^2\), disjoint ambient disks show that this net has at most \(Cn^6\) points, independently of the curve's length.  The reverse triangle inequality and the Lipschitz property of \(K_\triangle\) show that moving the center by \(d\) changes a scaled Gram average by at most \(Cd h^{-3}\), and changes the scaled outcome average by at most
\[
Cd h^{-3}\left(\frac1n\sum_{i=1}^n|Y_i|\right).
\]
The outcome envelope implies \(\sup_P\mathbb E_PY^2<\infty\).  A union bound over the net and the two arms, followed by the preceding oscillation inequalities, yields an event \(E_{n,h}\) with
\[
\sup_{P\in\mathcal P_{\mathrm{rect}}}P(E_{n,h}^c)\leq Cn^{-12}
\]
such that, uniformly over all boundary points, every denominator is at least \(\kappa_0/2\) and every centered numerator is at most
\[
C\sqrt{\frac{\log n}{nh^2}}.
\]
At \(h\asymp a_n\), the mesh oscillations are negligible because \(n^{-3}h^{-3}=o(a_n)\).  This argument uses an ambient-square net and hence requires no uniform boundary-length envelope or density-continuity modulus.

On the support of \(W_{i,t,x}(h)\), the H\"older assumption with \(q\geq1\) gives
\(|\mu_{t,P}(X_i)-\mu_{t,P}(x)|\leq Lh\).  Dividing the weighted decomposition by the denominator on \(E_{n,h}\) gives
\[
\sup_{x\in\mathcal B_P}|\widehat\mu_{t,h}(x)-\mu_{t,P}(x)|
\leq C\left\{h+\sqrt{\frac{\log n}{nh^2}}\right\}.
\]
On \(E_{n,h}^c\), a nonzero local-constant estimate is a convex combination of observed outcomes and hence is bounded by \(\max_i|Y_i|\); the zero-denominator branch is zero, and \(\sup_x|\tau_P(x)|\leq2L\).  Uniform sub-Gaussianity gives \(\sup_P\mathbb E_P\max_i|Y_i|^2\leq C\log n\).  Cauchy--Schwarz therefore makes the exceptional contribution \(O(n^{-6}\sqrt{\log n})\).  Taking the arm difference and then \(h\asymp a_n\) yields
\[
\sup_{P\in\mathcal P_{\mathrm{rect}}}\mathbb E_P\sup_{x\in\mathcal B_P}
|\widehat\tau^{\mathrm{LC}}_{n,h}(x)-\tau_P(x)|
\leq Ca_n,
\]
because \(h\) balances \(h\) with \(\sqrt{\log n/(nh^2)}\).  This proves the minimax upper bound with one data-driven member of the declared decision class.

For the converse, use the hypercube of \(M=M_n\) laws and centers from the angular-density packing.  Poissonize the observations with mean \(2n\) and attach independent continuous ordering marks.  Conditional on independent uniform hypercube bits, decompose the marked Poisson process into the disjoint bump cells \(Z_1,\ldots,Z_M\) and the complement \(Z_0\).  The channels are conditionally independent, \(Z_j\) depends only on bit \(j\), and \(Z_0\) is bit-independent.  Let \(S_j\) retain from cell \(j\) only outcomes, signed radii about \(x_j\), and ordering marks.  A coordinatewise distance decoder at \(x_j\), even after being given \(Z_{-j}\) and \(Z_0\), is measurable with respect to \((S_j,Z_{-j},Z_0)\).

Poisson relative entropy equals intensity times one-observation relative entropy.  The fixed-sample packing bound and equality of adjacent radial laws outside the local cell give
\[
\operatorname{KL}\{\mathcal L(S_j\mid\Omega_j=0),
\mathcal L(S_j\mid\Omega_j=1)\}
\leq2\alpha\log M.
\]
Since \(2\alpha<1\), the coordinatewise-overlap direct-product lemma implies probability at least \(1/2-o(1)\) of at least one bit error.  Midpoint thresholding of the estimated target at each center converts sup-norm error below \(c_1a_n/2\) into simultaneous correct decoding.  Hence every coordinatewise distance rule in the mean-\(2n\) Poisson experiment has Bayes expected sup-norm loss at least \(ca_n\).

To return to exactly \(n\) observations, start with any fixed-\(n\) rule.  In the marked Poisson experiment, if the count is at least \(n\), apply it to the \(n\) observations with smallest marks; otherwise output zero.  Conditional on having at least \(n\) points, the selected sample is i.i.d. from \(P\), while the Poisson lower-tail Chernoff inequality gives \(\Pr\{\operatorname{Poi}(2n)<n\}\leq e^{-cn}\).  Since \(\sup_x|\tau_P(x)|\leq2L\),
\[
ca_n\leq \sup_P\mathbb E_P^{\operatorname{Poi}(2n)}[\text{sup-loss}]
\leq \sup_P\mathbb E_P^n[\text{sup-loss}]+2Le^{-cn}.
\]
Taking the infimum over fixed-sample rules proves \(R_n^{\mathrm{dist}}(\mathcal P_{\mathrm{rect}})\geq c'a_n\).  Combining the two inequalities gives the asserted positive finite liminf and limsup constants.  The hard family is contained directly in the project-defined class; no relation to the different class of Cattaneo, Titiunik, and Yu, Theorem 6, is used.